\scanpage{13}
% archive sheet number: 12
1. Données de ramification sur un schéma.

Soit $S$ un schéma (préschéma dans l'ancienne terminologie), $\underline{\mathrm{Et}}_S$ le site étale de $S$. Nous appelerons \uline{donnée de ramification sur} $S$ \deleted{la donnée} un couple $\inserted{\underline{R}=}(\underline{C},\underline{s})$, où $\underline{C}$ est une catégorie fibrée sur $\underline{\mathrm{Et}}_S$ (\deleted{en \unsure{particulier}}\inserted{i.e.} c'est une catégorie $\underline{C}$ munie d'un foncteur fibrant $p:\underline{C}\quad\underline{\mathrm{Et}}_S$) et où $\underline{s}:\underline{C}\quad\underline{\mathrm{Fl}}_S$ est un foncteur cartésien, où $\underline{\mathrm{Fl}}_S$ désigne la catégorie fibrée des \deleted{f\unsure{ibreux|ibrés}} schémas relatifs sur les objets $S'$ de $\underline{\mathrm{Et}}_S$. On fera attention qu'on ne suppose pas que la catégorie fibrée $\underline{C}$ soit un champ (i.e. qu'on puisse recoller morphismes et objects de $\underline{C}$). Un objet $\inserted{M}$ de $\deleted{\underline{C}(S')}$ la catégorie fibre $\underline{C}_{S'}$ ($S'\quad\mathrm{Ob}\,\underline{\mathrm{Et}}_S$) sera appelé un modèle de ramification pour $\underline{R}$, et le schéma $\underline{s}(M)$ sur $S'$ sera appelé la réalisation géométrique de ce modèle. On fera attention qu'on ne suppose pas le foncteur $\underline{s}$ fidèle ni pleinement fidèle, donc les morphismes entre modèles de ramification ne peuvent être identifiés aux morphismes entre leurs réalisations géométriques. Dansles exemples envisagés par la suite, les $\underline{s}(M)$ sont finis sur $S'$.

Considérons le foncteur composé $\underline{C}\overset{\underline{s}}{\qquad}\underline{\mathrm{Fl}}_S\overset{\mathrm{source}}{\qquad}(\mathrm{Sch})$, d'où un foncteur (le foncteur ``réalisation géométrique'')
\[
\underline{C}\qquad(\mathrm{Sch})\,.
\]
Considérons d'autre part la catégorie fibrée sur$(\mathrm{Sch})$ des schémas relatifs étales et finis (dont la catégorie fibre en $T$ est donc la catégorie des revêtements étales finis de $T$), et soit $\underline{D}$ la catégorie fibrée sur $\underline{C}$ image inverse de la catégorie fibrée précédente par le foncteur réal'sation géométrique. La fibre de cette catégorie fibrée en l'objet $M$ de $\underline{C}$ est donc la catégorie des revêtements étales finis de la réalisation géométrique $\underline{s}(M)$ de $M$. O\textsuperscript{\unsure{N|n}} appelera \uline{revêtement de} \deleted{ty\unsure{p}}
% Physical right edge: only ty and a damaged p-like descending fragment of the cancelled word survive; no terminal e is visible.
